\documentclass[11pt]{article}
\usepackage[margin=2.5cm]{geometry}
\usepackage[spanish]{babel}
\usepackage[utf8]{inputenc}
\usepackage{amsmath}
\usepackage{listings}
\usepackage{xcolor}
\usepackage{hyperref}
\usepackage{booktabs}
\usepackage{parskip}

\lstset{
    basicstyle=\ttfamily\small,
    breaklines=true,
    frame=single,
    backgroundcolor=\color{gray!10},
    language=C
}

\title{Simulación de Ecosistema con OpenMP}
\author{Javier España, José Mérida, Angel Esquit \\ Computación Paralela y Distribuida}
\date{\today}

\begin{document}
\maketitle

Este proyecto implementa una simulación de ecosistema sobre una grilla 2D
(plantas, herbívoros, carnívoros) que avanza en \emph{ticks} discretos,
paralelizada con OpenMP. 

Repositorio: \url{https://github.com/AngelEsquit/MiniPRY-CPD}

\section{Reglas del Ecosistema}

Cada celda de la grilla guarda:

\begin{lstlisting}[language=c]
typedef struct {
    int type;    // EMPTY, PLANT, HERBIVORE, CARNIVORE
    int energy;
    int starve;  // ticks consecutivos sin comer
    int age;
} Cell;
\end{lstlisting}

\textbf{Plantas.} Cada tick, con 30\% de probabilidad, una planta se
expande a una celda vacía adyacente al azar; la planta original no
desaparece, sigue en su celda y crece un brote nuevo en la vecina. Muere
únicamente si es consumida por un herbívoro.

\textbf{Herbívoros.} Cada tick, un herbívoro pondera sus hasta 4 celdas
vecinas transitables (vacías, o con una planta que puede comer) como
suma ponderada de:

\begin{itemize}
    \item[+] $5 \times$ distancia de ese vecino al carnívoro más cercano
    \item[--] $2 \times$ distancia de ese vecino a la planta más cercana
    \item[+] $20$ si ese vecino tiene una planta (bono fijo por comer)
\end{itemize}

y se mueve al vecino de mayor puntaje (se queda quieto si ninguno es
transitable). Es un trade-off continuo entre huir del carnívoro más
cercano y acercarse a la planta más cercana, así podemos emular de manera
más cercana un comportamiento adecuado. Al comer gana energía, si llega a
$\geq 2\times$ su energía inicial, puede reproducirse (deja una cría en la celda
que acaba de desocupar). Muere por no comer, por vejez, o por ser comido.

\textbf{Carnívoros.} Misma estructura, pero sin comportamiento de huida.
Si hay un herbívoro adyacente lo cazan y si no, caminan
hacia el herbívoro más cercano. Si no hay ninguno cerca, deambulan. La
reproducción y la muerte siguen la misma lógica que los herbívoros, con sus
propias constantes de energía y hambre.

\textbf{Competencia entre especies.} Las tres especies deciden su
movimiento sobre la misma foto congelada de la grilla (nadie ve el
resultado de nadie más dentro del mismo tick), así que puede pasar que
varias celdas de especies distintas propongan la misma celda destino.
Por ejemplo, una planta expandiéndose a una celda vacía y un herbívoro
deambulando hacia esa misma celda. Eso se resuelve con una prioridad fija
por tipo de movimiento (más detalle en la Sección \ref{sec:resolve}):
cazar (20) $>$ huir (15) $>$ comer planta (10) $>$ deambular/expandirse
(5 / 1).

\section{Funcionamiento}

\subsection{Organización del Código}

El código está separado por responsabilidad en
\texttt{src/}:

\begin{lstlisting}[language=bash]
src/
  internal.h     Estado y helpers compartidos (privado)
  main.c          CLI
  core/           Estado, mapas de distancia, orquestacion del tick, apply
  species/        Una decision de movimiento por especie
  sync/           Resolucion de conflictos de destino
  io/             Salida de consola
  util/           RNG
\end{lstlisting}

\subsection{Estructuras Principales}

\begin{itemize}
    \item \texttt{Cell}: el estado de una celda (\texttt{type},
    \texttt{energy}, \texttt{starve}, \texttt{age}).
    \item \texttt{Move}: la propuesta de movimiento de una celda origen
    (\texttt{dst\_i}, \texttt{dst\_j}, \texttt{score}, \texttt{valid}),
    guardada en el arreglo \texttt{moves[]}.
    \item \texttt{claim\_score[]} / \texttt{claim\_src[]}: por cada celda
    destino, el mejor puntaje visto hasta ahora y quién lo propuso. Es el
    resultado de la resolución de conflictos.
    \item \texttt{moved[]} / \texttt{eaten[]}: banderas por celda origen
    para dos casos especiales (alguien que se movió de verdad, alguien
    que fue cazado en su propia celda) que la resolución por destino, por
    sí sola, no cubre.
    \item \texttt{Birth}: una cría en cola, con su celda y su tipo
    (\texttt{births[]}, \texttt{birth\_count}).
    \item \texttt{dist\_to\_plant[]}, \texttt{dist\_to\_herbivore[]},
    \texttt{dist\_to\_carnivore[]}: los mapas de distancia usados por
    herbívoros y carnívoros para decidir hacia dónde moverse.
\end{itemize}

\subsection{Ciclo del Tick}

Todo tick lee de \texttt{grid\_cur} y escribe en \texttt{grid\_next},
intercambiando los punteros recién al final (doble buffer). Si todos los
hilos leyeran y escribieran la misma grilla a la vez, el resultado
dependería del orden de ejecución de los hilos, que en OpenMP no está
definido. Con doble buffer, cada ronda paralela lee de un array que nadie
escribe y escribe en uno que nadie lee, así que no hace falta ningún
lock. El mismo truco se usa dentro de cada mapa de distancia (Sección
\ref{sec:distmaps}).

\begin{enumerate}
    \item Se recalculan tres mapas de distancia (a la planta,
    herbívoro y carnívoro más cercanos) desde \texttt{grid\_cur}.
    \item Cada especie propone, en paralelo, un movimiento por celda
    hacia \texttt{moves[]}.
    \item Se resuelven los conflictos de destino, llenando
    \texttt{claim\_src[]}, \texttt{moved[]} y \texttt{eaten[]}.
    \item Se aplica el resultado a \texttt{grid\_next}, usando esas tres
    estructuras.
    \item Se intercambian \texttt{grid\_cur} y \texttt{grid\_next} y se
    repite.
\end{enumerate}

\subsection{Mapas de Distancia}
\label{sec:distmaps}

Calcular "la planta más cercana" para cada celda es un BFS multi-fuente.
En vez de una cola compartida (su propio problema de sincronización), se
calcula por relajación en rondas, con el mismo truco de doble buffer:

\begin{lstlisting}[language=c]
buf_a[k] = (grid_cur[k].type == seed_type) ? 0 : DIST_INF;
for (round = 0; round < MAX_RELAX_RADIUS; round++) {
    #pragma omp parallel for collapse(2)
    for (i, j) {
        best = buf_a[IDX(i,j)];
        for cada vecino: best = min(best, buf_a[vecino] + 1);
        buf_b[IDX(i,j)] = best;
    }
    swap(buf_a, buf_b);
}
\end{lstlisting}

Por ronda, cada hilo procesa una celda: lee su propia distancia y la de
sus 4 vecinos en \texttt{buf\_a}, y escribe un único valor en
\texttt{buf\_b}. Como todos leen de \texttt{buf\_a} y nadie lo escribe
esa ronda, no hace falta coordinación entre hilos hasta el barrier
implícito al final del \texttt{parallel for}, que marca cuándo la
siguiente ronda puede arrancar.

Las rondas se acotan a \texttt{MAX\_RELAX\_RADIUS = 8}: a un organismo
solo le importa la dirección hacia el objetivo, no la distancia exacta en
toda la grilla.

\subsection{Lista de Movimientos}

Cada celda origen escribe su movimiento propuesto en
\texttt{moves[IDX(i,j)]}, su propio índice y ningún otro. Como ningún par
de hilos escribe jamás la misma posición, esta etapa no necesita
sincronización.

\subsection{Resolución de Conflictos}
\label{sec:resolve}

El único conflicto real es cuando dos celdas origen (de cualquier
especie) proponen la misma celda destino. Se resuelve con un
\texttt{omp\_lock\_t} por celda de la grilla:

\begin{lstlisting}[language=c]
omp_set_lock(&cell_locks[dst]);
if (moves[k].score > claim_score[dst]) {
    claim_score[dst] = moves[k].score;
    claim_src[dst]   = k;
}
omp_unset_lock(&cell_locks[dst]);
\end{lstlisting}

Cada lock protege una sola celda, así que dos hilos yendo a destinos
distintos nunca se bloquean entre sí.

La prioridad por score cubre exclusivamente empates sobre una misma
celda destino; por las reglas de movimiento, la celda de un herbívoro
vivo solo puede ser destino de un carnívoro cazándolo, así que ahí el
único empate posible es entre dos carnívoros cazando al mismo
herbívoro.

Después de resolver corren dos pasadas cortas de etiquetado, sin locks
(cada una escribe un índice distinto por iteración, libre de carrera):
\texttt{moved[]} marca qué celdas se vaciaron de verdad, para no
duplicar a alguien que se movió como si además se hubiera quedado
quieto; \texttt{eaten[]} marca herbívoros cazados en su propia celda de
origen, cubriendo el caso en el que ese mismo herbívoro también ganó,
en una celda destino distinta, su propio movimiento hacia otra parte.
\texttt{core/apply.c} revisa \texttt{eaten[]} antes de colocar a
cualquier ganador de movimiento.

\subsection{Nacimientos}

Una cría se coloca en la celda que su padre acaba de desocupar, una
celda que también se procesa, en la misma pasada paralela, como destino
de otra celda cualquiera. Escribirla directamente ahí sería una carrera,
así que cada nacimiento reserva un lugar en una cola con un contador
atómico (\texttt{\#pragma omp atomic capture}) y se coloca en una pasada
aparte, después de que la pasada principal termina.

\subsection{Números Aleatorios}

Cada hilo tiene su propia semilla (\texttt{seeds[tid]}) y usa
\texttt{rand\_r}, reentrante, para no competir por el estado global de
\texttt{rand()}.

\section{Compilación y Ejecución}

\begin{lstlisting}[language=bash]
make            # genera build/*.o y el binario ./ecosystem
make clean      # borra build/ y el binario
./ecosystem <filas> <cols> <ticks> <hilos>
./ecosystem 60 60 100 8
\end{lstlisting}

La población inicial se deriva del tamaño de la grilla: 30\% plantas,
10\% herbívoros, 5\% carnívoros. Cada 10 ticks (y en los ticks 1 y 5) se
imprime el conteo de población por especie y, si la grilla es de 30x30 o
menos, un mapa de caracteres de la distribución.

\section{Resultados}

En corridas de prueba (60x60, 40 ticks) se observa el patrón esperado de
un sistema depredador-presa: plantas fluctuando con la presión de los
herbívoros, herbívoros creciendo con la comida disponible y cayendo con
la caza, carnívoros decayendo si la población de herbívoros colapsa. No
se observó duplicación de organismos ni poblaciones que excedan la
capacidad de la grilla en ninguna corrida.

Tiempo total de simulación variando el número de hilos (grilla 200x200,
50 ticks, máquina con 20 núcleos lógicos):

\begin{center}
\begin{tabular}{@{}rrrr@{}}
\toprule
Hilos & Tiempo total (s) & Tiempo/tick (ms) & Speedup \\
\midrule
1 & 0.2138 & 4.28 & 1.00$\times$ \\
2 & 0.1297 & 2.59 & 1.65$\times$ \\
4 & 0.0878 & 1.76 & 2.44$\times$ \\
8 & 0.0860 & 1.72 & 2.49$\times$ \\
\bottomrule
\end{tabular}
\end{center}

El speedup se aplana entre 4 y 8 hilos: cada tick lanza más de veinte
regiones paralelas distintas (8 rondas $\times$ 3 mapas de distancia, más
decisión/resolución/aplicación), y cada una paga un costo fijo de
arranque de hilos que, en una grilla de este tamaño, empieza a pesar más
que el trabajo ganado por hilo adicional. Para grillas más grandes se
esperaría que ese punto se corra hacia más hilos.

\end{document}
